% hoeffding-two-perspectives.tex
% Notes on Høffding's internal/external distinction, for vtfys Week 5

\documentclass[12pt, a4paper]{article}

\usepackage[utf8]{inputenc}
\usepackage[T1]{fontenc}
\usepackage{lmodern}
\usepackage{microtype}
\usepackage[top=3cm, bottom=3cm, left=3.5cm, right=3.5cm]{geometry}
\usepackage[backend=biber, style=authortitle-ibid, ibidtracker=strict]{biblatex}

\usepackage{csquotes}
\usepackage{hyperref}
\hypersetup{colorlinks=true, linkcolor=blue, citecolor=blue, urlcolor=blue}

\title{\vspace{-1.5em}H{\o}ffding on the Two Perspectives:\\
Inner and Outer as Irreconcilable\vspace{-0.5em}}
\author{Notes for vtfys Week 5}
\date{}

\begin{document}

\maketitle

\subsection*{Background}

Harald H{\o}ffding's \emph{Outlines of Psychology} (1882; English
translation 1891) develops a sustained account of the relationship
between psychological description and physiological description.
H{\o}ffding holds a version of psychophysical parallelism---mind and
body are not two substances in interaction, but one and the same
reality seen from two different sides. The parallel between conscious
processes and brain processes is real and deep. Yet the two
perspectives are not reducible to each other: they cannot be
translated into a common language, and no amount of external
observation yields the inner perspective. The following passages
document this claim with precision.

\subsection*{The Two Sides}

H{\o}ffding introduces the parallelism in Chapter~II (\emph{Mind
and Body}) using Fechner's simile:

\begin{quote}
Physiology and psychology deal with the same matter seen from two
different sides, and there can no more be dispute between them than
between the observer of the convex and the observer of the concave
side of a curve.\footnote{H{\o}ffding, \emph{Outlines of Psychology},
trans.\ M.\ E.\ Lowndes (London: Macmillan, 1891), p.~70.}
\end{quote}

\noindent
The same chapter establishes that the two perspectives cannot be
collapsed into one. Writing in the tradition of Spinoza and
Leibniz---both of whom insisted that thoughts must be explained by
thoughts and movements by movements---H{\o}ffding notes that the two
orders of description ``cannot be reduced to a common
measure.''\footnote{Ibid., p.~69.} This is not merely an epistemic
limitation; it is a structural feature of the two descriptions.

\subsection*{Two Languages with No Common Original}

The most striking formulation of the irreducibility appears a few
pages earlier in the same chapter. H{\o}ffding has been arguing that
for every conscious phenomenon there is a corresponding material
process (and vice versa), and that the two are not causally related
but constitute two manifestations of the same underlying reality.
He then writes:

\begin{quote}
It is as though the same thing were said in two
languages.\footnote{Ibid., p.~65.}
\end{quote}

\noindent
The metaphor is immediately qualified in a way that sharpens its
force:

\begin{quote}
The two languages, in which the same thought is here expressed, we
are not able to trace back to a common original
language.\footnote{Ibid., p.~66.}
\end{quote}

\noindent
This is the crux. H{\o}ffding does not deny the unity of the
underlying reality. But he holds that the two descriptive perspectives
share no common original vocabulary into which both could be
translated. The inner and the outer are irreconcilable not because
they describe different things, but because there is no third language
in which the description could be given.

\subsection*{The Inner as the Only Starting Point}

H{\o}ffding draws out the epistemological consequence in Chapter~I
(\emph{Subject and Method}). Conscious life is accessible only from
the inside:

\begin{quote}
If we wish to gain a knowledge of conscious life, we must study it,
above all, where it is directly accessible to us---namely, in our own
consciousness. This immediate experience is also the only source
whence the physiologist can determine the significance for mental
life of the various organs of the brain.\footnote{Ibid., p.~11.}
\end{quote}

\noindent
The external description---the physiologist's description of the
brain---cannot itself yield the inner content it is correlated with.
Only first-person access provides that. H{\o}ffding puts the
point even more sharply in the context of ruling out both spiritualism
and materialism as starting points for psychology:

\begin{quote}
Psychological experience gives only the internal mental phenomena
themselves---not the manner in which they are connected with other
phenomena.\footnote{Ibid., p.~14.}
\end{quote}

\subsection*{The Leibniz Connection}

H{\o}ffding's position is structurally identical to the conclusion
of Leibniz's mill argument (\emph{Monadologie}, \S{}17), though he
arrives at it by a different route. Leibniz imagined the brain
enlarged to the size of a mill, so that one could walk around inside
and observe all its physical workings---and argued that one would
never, in this way, encounter \emph{perception} or \emph{thought}.
H{\o}ffding, working from psychophysical parallelism rather than
from Leibniz's metaphysics, reaches the same conclusion: the outer
perspective, however refined, cannot generate the inner one. He
explicitly cites the \emph{Monadologie} elsewhere in the
\emph{Outlines},\footnote{Ibid., p.~46 n.\ 1, citing
\emph{Monadologie} \S{}24.} and his \emph{History of Modern
Philosophy} contains a detailed treatment of Leibniz.

The significance for Bohr is direct. When Bohr insists that the
experimental arrangement must be described in classical terms---that
the classical description is a condition of unambiguous communication,
not a contingent choice---he is applying the H{\o}ffding-style
principle to quantum mechanics: there is a perspective from which
unambiguous statements can be made, and it cannot be replaced by the
quantum description of the apparatus itself, any more than the inner
perspective can be replaced by the outer one. Complementarity, on
this reading, is the quantum-mechanical instantiation of the
principle that the two languages share no common original.

\end{document}

%%% Local Variables:
%%% mode: latex
%%% TeX-master: t
%%% End:
